% MT017 / explanation — canonical editable content.
线性边界满足\(w^Tx+b=0\)：二维时是一条直线，更高维对应一个超平面。非线性边界不能只用这个形式表达，例如圆弧或分段折线。题目问同一种模型换数据后，能否形成两类形状，而不是问它是否总能找到任意边界。



\begin{itemize}

\item \textbf{A选。} Gaussian NB比较两类的高斯分数。若两类共享每一维的方差，展开后\(x_j^2\)项抵消，只剩一次项；方差不同时，二次项通常保留，所以也能产生非线性边界。完整展开见\StudyRef{Answers}{mt059}{MT059}。

\item \textbf{B按原解选。} RBF核根据到训练样本的距离形成局部响应，组合后能产生弯曲边界；对称数据及系数也可以得到直线分隔。使用RBF不等于边界必须弯曲。

\item \textbf{C按原解选。} 决策树桩只做一次阈值判断，如\(x_1>t\)。多个树桩加权投票，可只使用一个有效阈值形成直线，也可结合不同阈值形成分段边界。有限树桩不保证精确表示任意光滑曲线。

\item \textbf{D不选。} 注意核是\((x^Tx')^2\)，没有“\(+1\)”。它只提供二次乘积特征，例如\(x_1^2,x_1x_2,x_2^2\)，缺少一般的一次项。把输入从\(x\)换成\(-x\)，这些二次乘积不变，分数也不变；它无法像一般线性分类器那样，按一个方向的正负区分这两个点。

\item \textbf{E不选。} 改变正则强度\(\alpha\)只改变权重大小，分数仍是\(w^Tx+b\)。若没有另加非线性特征，就不会因此获得曲线边界。

\end{itemize}


选读：某些二次式的零集合能退化成直线，但这不等于它能像一般线性分类器那样给直线两侧分配不同类别。


\textbf{出处：} 2021A期中 Q4，原卷第2页（5分）


\textbf{答案依据：} 2021A配套解答 / 2021A配套解答副本 Q4，随卷答案ABC。
